% ============================================================
% 分片：s02c  对应原文 PDF 第 10--14 页
% 负责章节：2.1.4 Equivariant cohomology（余下部分）；2.2 Mirror symmetry from the worldsheet；2.2.1；2.2.2
% 译者：translator-02   审校状态：初稿
% ============================================================

% [跨页接续: 原文第 9 页末句 "...it will if we restrict our attention to $J$-invariant operators." 按 SPEC §3 语义完整取于 s02b；本分片自 "See [27] ..." 起。]

% ==== 原文 PDF 第 10 页 ====

见 [27] 关于这一主题面向物理学家的详尽介绍。为简单起见，我们只讨论阿贝尔情形。

这一想法基于等变上同调（equivariant cohomology）的 Cartan 模型：在带有（比如说）由矢量场 $V$ 生成的 $U(1)$ 作用的光滑流形 $M$ 上，de Rham 微分 $d_M$ 有自然的推广，它作用在微分形式与多项式代数作张量积所得的对象 $\mathbb{C}[\sigma]\otimes\Omega^\bullet(M)$ 上\footnote{更一般地，若不专门讨论阿贝尔情形，则得到 $S(\mathfrak{g}^*)\otimes\Omega^\bullet(M)$，即取值于 Lie 代数对偶的对称代数中的形式。}，由 $d_\sigma=d_M+\sigma\iota_V$ 给出，其中 $\iota_V$ 表示与矢量场 $V$ 的缩并，$\sigma$ 是称为等变参数（equivariant parameter）的形式变量。注意还有 $d_\sigma\sigma=0$。Cartan 公式（关于 $V$ 的 Lie 导数 $\mathcal{L}_V$ 可表为 $\mathcal{L}_V=d_M\iota_V+\iota_V d_M$）表明 $d_\sigma^2=\sigma\mathcal{L}_V$。阿贝尔情形的等变上同调本质上就是 $U(1)$ 不变微分形式的 $d_\sigma$ 上同调，而 $U(1)$ 不变微分形式即满足 $\mathcal{L}_V\omega=0$ 的形式 $\omega\in\Omega^\bullet(M)$。更精确地说，我们在这些微分形式上再添上形式参数 $\sigma$，由此实现 $M$ 的 $U(1)$ 等变上同调：

\begin{equation}
H^\bullet_{U(1)}(M) = H\big((\mathbb{C}[\sigma]\otimes\Omega^\bullet(M))^{U(1)}, d_\sigma\big) = H\big(\Omega^\bullet(M)^{U(1)}[\sigma], d_\sigma\big).
\tag{2.8}
\end{equation}
% [Tutor Note: $d_\sigma^2=\sigma\mathcal{L}_V$ 的来历：由 Cartan 公式 $\mathcal{L}_V=d_M\iota_V+\iota_V d_M$ 以及 $d_M^2=0$、$\iota_V^2=0$（$\sigma$ 为偶元素，与 $d_M$、$\iota_V$ 对易），直接算得 $d_\sigma^2=(d_M+\sigma\iota_V)^2=d_M^2+\sigma(d_M\iota_V+\iota_V d_M)+\sigma^2\iota_V^2=\sigma\mathcal{L}_V$。因此 $d_\sigma$ 并非处处平方为零；它只在 $\mathcal{L}_V\omega=0$ 的不变形式上才平方为零。又因等变参数 $\sigma$ 是自由的，$d_\sigma^2\omega=0$ 与 $\mathcal{L}_V\omega=0$ 等价，故可在系数环中约去 $\sigma$，得到 (2.8) 第二个等号所写的 $H(\Omega^\bullet(M)^{U(1)}[\sigma],d_\sigma)$。]
% [Tutor Note: (2.8) 所写的即等变上同调的 Cartan 模型：把 $H^\bullet_G(M)$ 实现为 $(S(\mathfrak{g}^*)\otimes\Omega^\bullet(M))^G$ 上的上同调。阿贝尔情形下 $S(\mathfrak{g}^*)=\mathbb{C}[\sigma]$，即"微分形式与多项式的张量积、取不变、再用 $d_\sigma$ 取上同调"；脚注给出非阿贝尔时的对应写法，其中 $S(\mathfrak{g}^*)$ 是 Lie 代数对偶的对称代数。]

若给形式参数 $\sigma$ 以上同调度 2，则等变上同调具有 $\mathbb{Z}$ 分次，总度数为形式度数加上 $\sigma$ 中次数的两倍。这一修改在研究如下情形时经常出现，例如：1）带中心荷的扭变理论；2）存在 $\Omega$-背景（Omega-background）[28] 的情形，其中 $M$ 是适当的场空间，$U(1)$ 对称性表现为时空中绕某轴的转动；或 3）研究某些超引力背景中的超对称理论。这些考虑的一种推广也是超对称局域化（supersymmetric localization）的基础，我们在附录 A 中简略勾画。
% [Tutor Note: 要求 $\sigma$ 的上同调度为 2 是为了让 $d_\sigma$ 提高总度数 1：$d_M$ 提高形式度 1 而不改变 $\sigma$ 的次数，$\sigma\iota_V$ 中 $\sigma$ 提高 2 次、缩并 $\iota_V$ 降低形式度 1，两者之和仍为 1。若把 $\sigma$ 的度数取成 1，$d_\sigma$ 将是奇算子但不提高度数，等变上同调就不再按通常意义分次。核对局域化公式（附录 A）中 $\sigma$ 幂次的计数时须牢记这一度数约定。]
% [背景补充: $\Omega$-背景是 Nekrasov 引入的一类超对称背景形变：在时空中取两个相互对易的 $U(1)$ 转动，使路径积分局域化到其固定点集上，从而得到只依赖固定点邻域数据的配分函数；此处 $M$ 取为场空间，作用即在 $M$ 上，故正落在等变上同调的框架内。]

最后，我们简要提及 Costello 与 Li [11] 提出的直接扭变超引力理论的一个方案：它以经典超引力数据具体地修改我们概述的扭变程序——一个时空流形 $\Sigma$，以及 $\Sigma$ 上带联络的丛 $P_{\mathrm{Spin}}$、$P_R$。

\begin{enumerate}
\item 不选取幂零的 $Q$，而是让与超平移相联系的玻色鬼场取真空期望值。
\item 扭变同态 $\iota$ 被替换为时空上 $G$-丛的选取 $P_G\to\Sigma$，其中还包含联络的选取，使得丛 $P_{\mathrm{Spin}}\times P_R$ 由同态 $G\to Spin(d)\times G_R$ 诱导得到。
\item 若适用，映射 $\gamma:U(1)\to G_R$ 的类比是：在 $G_R$ 中选取一个平凡的 $U(1)$ 子丛，并使联络在其上限制为零。
\end{enumerate}
% [Tutor Note: 与前面在平坦空间上选 $\iota$ 的做法相比，这一方案把扭变"丛化"：$\iota$ 是 $Spin(d)\to G_R$ 的群同态，而这里取 $G$-丛 $P_G$ 及同态 $G\to Spin(d)\times G_R$，把 $P_{\mathrm{Spin}}\times P_R$ 实现为 $P_G$ 的伴丛。三条修改分别对应：以超平移鬼场的 VEV 代替幂零超荷、以 $G$-丛代替同态 $\iota$、以 $G_R$ 中的平凡 $U(1)$ 子丛代替等变参数 $\gamma$。其中"联络在子丛上限制为零"正是等变上同调中 $\sigma$ 与 $G_R$ 生成元配对的具体几何实现。]

\subsection{从世界面看镜像对称}

镜像对称，以及包含膜（brane）范畴在内的其各种丰富形式，也许是对数学影响最为深远的弦对偶。至少，它一直是弦论与数学互动史上最突出的历史范例之一。
% ==== 原文 PDF 第 11 页 ====
镜像对称的扭变版本（包含膜范畴）常被称为同调镜像对称（homological mirror symmetry），但我们强调：镜像对称是完整物理（未扭变）理论之间的对偶。

镜像对称是 IIA 型与 IIB 型弦论在两张不同 Calabi--Yau 流形上的物理等价，这两张流形称为镜像流形（mirror manifolds）。我们可以微扰地研究对偶的每一边，例如使用世界面形式。此外，我们可以扭变超弦世界面，以触及完整物理对偶的一个数学上丰富但远为简单的子集；在该子集中，既能对偶性作严格的检验，也能产生新的数学结果。

这里，我们说明弦论中镜像对称研究（以及与之相关的二维 $\mathcal{N}=(2,2)$ QFT 研究）所依赖的一些基本要素，为第 2.3 节中关于同调镜像对称的一些相对基本的陈述作准备。虽然我们聚焦于二维 $\mathcal{N}=(2,2)$ QFT，我们特别关注 $U(1)_A\times U(1)_V$ $R$ 对称性，因为两个 $U(1)$ 的保持是在红外获得超共形不变性的必要条件，而后者又是弦世界面理论的必要条件。如上一节所述，这也意味着此类理论容许两种（至少 $\mathbb{Z}$ 分次的）拓扑扭变，它们将在第 2.3 节中占据中心地位。

\subsubsection{二维 $\mathcal{N}=(2,2)$ 超代数的镜像对称}

无须深入任何具体理论与拉氏量，我们就已能描述镜像对称是什么。正如 T 对偶的情形，它来自一个看似无害、却在时空上后果剧烈的世界面同构，而这一点在 $\mathcal{N}=(2,2)$ 超代数层面就已可见。（这或许并不太令人惊讶，因为我们此前已讨论过镜像对称如何可视为某个 T 对偶序列。）

我们在第 2.1.1 节简要提及了二维 $\mathcal{N}=(2,2)$ 超对称代数的结构，现在稍详细地描述它。记 $J$ 为世界面转动的生成元，使得例如 $[J,P_z]=-P_z$。如上所述，二维 $\mathcal{N}=(2,2)$ 超对称代数有 4 个超荷 $Q_\pm,\bar Q_\pm$，其中下标 $\pm$ 表示超荷的自旋／手征性：

\begin{equation}
[J,Q_\pm] = \pm\tfrac12 Q_\pm, \qquad [J,\bar Q_\pm] = \pm\tfrac12 \bar Q_\pm.
\tag{2.9}
\end{equation}

除了上面给出的反对易关系外，我们还可以引入两类复中心荷 $Z$、$\tilde Z$ 及其复共轭 $Z^*$、$\tilde Z^*$；带有中心荷的二维 $\mathcal{N}=(2,2)$ 超对称代数由如下括号给出：

\begin{equation}
\begin{aligned}
[Q_+,\bar Q_+] &= 2P_{\bar z} & [Q_-,\bar Q_-] &= -2P_z\\
[\bar Q_+,\bar Q_-] &= Z & [Q_+,Q_-] &= Z^*\\
[\bar Q_+,Q_-] &= \tilde Z & [Q_+,\bar Q_-] &= \tilde Z^*
\end{aligned}\tag{2.10}
\end{equation}
% [存疑: 原文正文称这些为 anti-commutation relations，但 (2.9)--(2.11) 一律用方括号 $[\ ,\ ]$ 书写超荷之间的关系；超代数中 $[\ ,\ ]$ 通常表示分次对易子（对两个奇元素即为反对易子），故原文多半只是未明说该约定，并非笔误。译文照抄方括号，读者按分次对易子理解即可。]

然而，由于这些中心项必须与所有其他代数元素对易，若 $R_V$ 守恒，则 $Z$、$Z^*$ 必须为零；若 $R_A$ 守恒，则 $\tilde Z$、$\tilde Z^*$ 必须为零。尽管如此，引入复质量 $Z$（相应地扭复质量）会破坏 $B$ 型超荷
% ==== 原文 PDF 第 12 页 ====
$Q_B$（相应地 $A$ 型超荷 $Q_A$）的幂零性，如我们在第 2.1.4 节所见。\footnote{这一比较可以做得相当精确：二维 $\mathcal{N}=(2,2)$ 理论的复质量形变，通常来自开启背景矢量多重态的一个标量分量 $\sigma$，它与（某个环面上的）味对称性 $F$ 耦合，于是 $Z$ 在荷为 $n$ 的超选择扇区上作用为 $n\sigma$。二维 $\mathcal{N}=(2,2)$ 理论的扭复质量形变则常来自 Fayet--Iliopoulos 项，并由二维 $\theta$ 角复化。} 这些中心荷给出了各类感兴趣理论中 BPS 孤子的质量，因此 $Z$（相应地 $\tilde Z$）常被称为复质量（相应地扭复质量），不过本讲义不研究它们。

上述代数有若干外自同构，即不是由与某个固定元素的对易所诱导的代数对称性。通常有 $\mathbb{Z}_2$ 宇称变换，它把 $Q_\pm\to Q_\mp$、$\bar Q_\pm\to\bar Q_\mp$；还有 $\mathbb{Z}_2$ 电荷共轭，它把 $Q_\pm\leftrightarrow\bar Q_\pm$。但还有第三种 $\mathbb{Z}_2$ 镜像自同构，它只交换左移超荷：$Q_-\leftrightarrow\bar Q_-$。为保持 $R$ 对称群，我们看到镜像对称必须交换轴 $U(1)_A$ 与矢量 $U(1)_V$ $R$ 对称性：$R_A\leftrightarrow R_V$。类似地，镜像对称必须交换复质量 $Z$ 与扭复质量 $\tilde Z$。最后注意，镜像自同构保持全纯超荷 $Q_H=\bar Q_+$，但交换两个拓扑超荷 $Q_A := \bar Q_+ + Q_- \leftrightarrow Q_B = \bar Q_+ + \bar Q_-$。

我们说两个 $\mathcal{N}=(2,2)$ 理论 $\mathcal{T}$ 与 $\tilde{\mathcal{T}}$ 互为镜像，若这两个理论之间存在等价（例如：其物理 Hilbert 空间中的态、配分函数、局域与延展算子的关联函数等之间的一个对应），且该等价使超对称生成元按上述镜像对合相互交换。特别地，由于镜像自同构交换两个拓扑超荷，$\mathcal{T}$ 的 $A$ 型扭变必等价于 $\tilde{\mathcal{T}}$ 的 $B$ 型扭变，反之亦然：$\mathcal{T}^A\cong\tilde{\mathcal{T}}^B$，$\mathcal{T}^B\cong\tilde{\mathcal{T}}^A$。正如我们从这一非常一般的讨论所预期的，镜像对称并非仅限于 Calabi--Yau 的现象，尽管它正是在该背景下被发现的 [29]。
% [Tutor Note: 镜像自同构对 (2.10) 的作用可以逐项核对：它把 $Q_-\leftrightarrow\bar Q_-$、$R_A\leftrightarrow R_V$、$Z\leftrightarrow\tilde Z$，于是 $[Q_+,\bar Q_-]=\tilde Z^*$ 与 $[Q_+,Q_-]=Z^*$ 互换，$[\bar Q_+,Q_-]=\tilde Z$ 与 $[\bar Q_+,\bar Q_-]=Z$ 互换，而 $[Q_+,\bar Q_+]=2P_{\bar z}$、$[Q_-,\bar Q_-]=-2P_z$ 也随 $P_z\leftrightarrow P_{\bar z}$ 与符号约定一起互换。这给出"镜像对称是超代数的外自同构"这一说法的具体含义。]

\subsubsection{手征与扭手征超多重态的镜像对称}

上文我们看到，镜像对称可解释为二维 $\mathcal{N}=(2,2)$ 超对称代数的某种外自同构。本小节描述 $\mathcal{N}=(2,2)$ 超代数的若干简单表示，称为手征多重态（chiral multiplets）与扭手征多重态（twisted chiral multiplets）。尽管本讲义不会涉及，还存在矢量多重态（vector multiplets）与扭矢量多重态（twisted vector multiplets），用于构造超对称规范理论。有一些有用的规范理论类，称为规范线性 $\sigma$ 模型（gauged linear sigma models, GLSM），例如可以被设计成在红外流向以 Calabi--Yau 为靶空间的非线性 $\sigma$ 模型。
% [背景补充: GLSM 由 $U(1)^k$ 规范场与手征多重态构成，其红外不动点给出非线性 $\sigma$ 模型；Witten 用它把 Calabi--Yau 的凯勒模空间实现为 GLSM 的真空模空间，因而成为检验镜像对称的标准工具。]

我们将要关注的二维 $\mathcal{N}=(2,2)$ 超对称代数的表示，可以用二维 $\mathcal{N}=(2,2)$ 超空间来打包。除复玻色坐标 $z,\bar z$ 之外，我们引入四个旋量性的费米坐标 $\theta^\pm,\bar\theta^\pm$。二维 $\mathcal{N}=(2,2)$ 超对称代数可实现为超空间上的平移：引入费米矢量场

\begin{equation}
\begin{aligned}
Q_+ &= \partial_{\theta^+} + i\bar\theta^+\partial_{\bar z} & \bar Q_+ &= -\partial_{\bar\theta^+} - i\theta^+\partial_{\bar z}\\
Q_- &= \partial_{\theta^-} - i\bar\theta^-\partial_z & \bar Q_- &= -\partial_{\bar\theta^-} + i\theta^-\partial_z
\end{aligned}\tag{2.11}
\end{equation}

由此可得所需的反对易子，例如 $[Q_+,\bar Q_+]=-2i\partial_{\bar z}=2P_{\bar z}$。轴 $U(1)_A$ 与矢量 $U(1)_V$ $R$ 对称性自然地作为费米坐标的转动出现在超空间上：

% ==== 原文 PDF 第 13 页 ====
% [本页以上一句末尾 "coordinates:" 的续接开始；(2.12) 为第 13 页的首个公式。]

\begin{equation}
(e^{i\alpha},e^{i\beta})\in U(1)_A\times U(1)_V \rightsquigarrow (\theta^\pm,\bar\theta^\pm)\mapsto (e^{\mp i\alpha-i\beta}\theta^\pm,\ e^{\pm i\alpha+i\beta}\bar\theta^\pm).
\tag{2.12}
\end{equation}

这一实现给出二维 $\mathcal{N}=(2,2)$ 代数的若干自然表示，其载体是超空间上的函数（或更一般地，某个丛／层在超空间上的截面）$S(z,\bar z;\theta^\pm,\bar\theta^\pm)$，称为超场（superfield）。坐标 $\theta^\pm,\bar\theta^\pm$ 的费米性质意味着关于 $\theta^\pm=\bar\theta^\pm=0$ 的 Taylor 展开必然有限；确实，一般超场 $S(z,\bar z;\theta^\pm,\bar\theta^\pm)$ 包含 $2^4=16$ 个场。一旦选定超场的内禀矢量与轴 $R$ 荷 $q_V$、$q_A$，即费米 Taylor 展开中常数项的荷，各分量场的荷便由下式唯一确定：

\begin{equation}
\begin{aligned}
e^{i\alpha R_A}S(z,\bar z;\theta^\pm,\bar\theta^\pm) &= e^{i\alpha q_A}S(z,\bar z;e^{\mp i\alpha}\theta^\pm,e^{\pm i\alpha}\bar\theta^\pm)\\
e^{i\beta R_V}S(z,\bar z;\theta^\pm,\bar\theta^\pm) &= e^{i\beta q_V}S(z,\bar z;e^{-i\beta}\theta^\pm,e^{i\beta}\bar\theta^\pm)
\end{aligned}\tag{2.13}
\end{equation}
% [Tutor Note: (2.12) 与 (2.13) 的相位约定值得逐项核对：轴 $U(1)_A$ 的相位随 $\pm$ 指标变号——$\theta^\pm$ 得 $e^{\mp i\alpha}$、$\bar\theta^\pm$ 得 $e^{\pm i\alpha}$；矢量 $U(1)_V$ 的相位则与 $\pm$ 指标无关——所有 $\theta^\pm$ 都得 $e^{-i\beta}$、所有 $\bar\theta^\pm$ 都得 $e^{+i\beta}$。这一"A 轴敏感于手征指标、V 只区分 $\theta$ 与 $\bar\theta$"的差别正是 $A$ 型与 $B$ 型扭变分别用轴／矢量 $R$ 对称性的原因。若把 (2.13) 两行的相位抄混，后续分量场的 $R$ 荷与 (2.12) 的超空间转动就会不相容。]

类似地，组成场的费米宇称由超场 $S(z,\bar z;\theta^\pm,\bar\theta^\pm)$ 的内禀费米宇称决定。例如，若 $S=a+\theta^+\alpha+\ldots$ 是玻色（相应地费米）超场，则 $a$ 是玻色子（相应地费米子），$\alpha$ 是费米子（相应地玻色子）。

结果表明，一般超场 $S(z,\bar z;\theta^\pm,\bar\theta^\pm)$ 并不给出二维 $\mathcal{N}=(2,2)$ 超代数的不可约表示。超空间再次给我们帮助：它提供自然的微分算子 $D_\pm$、$\bar D_\pm$（称为超导数，superderivatives），它们与超对称代数反对易：

\begin{equation}
\begin{aligned}
D_+ &= \partial_{\theta^+} - i\bar\theta^+\partial_{\bar z} & \bar D_+ &= -\partial_{\bar\theta^+} + i\theta^+\partial_{\bar z}\\
D_- &= \partial_{\theta^-} + i\bar\theta^-\partial_z & \bar D_- &= -\partial_{\bar\theta^-} - i\theta^-\partial_z
\end{aligned}\tag{2.14}
\end{equation}

由于这些算子与 $Q_\pm$、$\bar Q_\pm$ 对易，可以用它们约束超场的分量。例如，手征超场（chiral superfield）$\Phi(z,\bar z;\theta^\pm,\bar\theta^\pm)$ 要求满足 $\bar D_\pm\Phi=0$，其复共轭反手征超场（anti-chiral superfield）$\bar\Phi$ 满足 $D_\pm\bar\Phi=0$。类似地，扭手征超场（twisted chiral superfield）$\tilde\Phi$ 满足 $\bar D_+\tilde\Phi=D_-\tilde\Phi=0$（对复共轭场有类似结论）。
% [存疑: 原文上一句称超导数"与超对称代数反对易"，此处却说"与 $Q_\pm$、$\bar Q_\pm$ 对易"；按超代数约定，二者应统一理解为分次对易子为零（即 $\{D_\pm,Q_\pm\}=\{D_\pm,\bar Q_\pm\}=0$），原文用词略有不一。译文照原文作"对易"，读者按分次对易子理解即可。]
% [Tutor Note: 逐项比较 (2.11) 与 (2.14)：$Q$ 与 $D$ 只差在含 $\theta\partial$ 的项上符号相反——例如 $Q_+$ 中为 $+i\bar\theta^+\partial_{\bar z}$，而 $D_+$ 中为 $-i\bar\theta^+\partial_{\bar z}$。这正是反对易关系 $\{Q_+,D_+\}=0$ 的显式体现，可由 $\{\partial_{\theta^+},\theta^+\}=1$、$\{\bar\theta^+,\bar\theta^+\}=0$ 逐项验证。特别提醒：(2.14) 的 $D_-$ 中 $i\bar\theta^-\partial_z$ 前是正号，而 (2.11) 的 $Q_-$ 中是负号，抄写时极易出错。]

引入移位坐标 $y=z+i\theta^-\bar\theta^-$、$\bar y=\bar z-i\theta^+\bar\theta^+$ 是方便的。手征性约束 $\bar D_\pm\Phi=0$ 意味着该超场只经由 $y$、$\bar y$ 与 $\theta^\pm$ 依赖各坐标。于是，可把手征超场写成\footnote{用移位坐标写出的手征超场表达式，展示了手征多重态的一种简洁性，而这种简洁性在其完整分量展开中可以部分地被掩盖。完整展开为
\[
\begin{aligned}
\Phi = {} & \phi + \theta^+\psi_+ + \theta^-\psi_- + \theta^+\theta^-F + \theta^+\bar\theta^+(-i\partial_{\bar z}\phi) + \theta^-\bar\theta^-(i\partial_z\phi)\\
& + \theta^+\theta^-\bar\theta^+(i\partial_{\bar z}\psi_-) + \theta^+\theta^-\bar\theta^-(i\partial_z\psi_+) + \theta^+\theta^-\bar\theta^+\bar\theta^-(-\partial_{\bar z}\partial_z\phi).
\end{aligned}
\]
}：

\begin{equation}
\Phi(y,\bar y;\theta^\pm,\bar\theta^\pm) = \phi(y,\bar y) + \theta^+\psi_+(y,\bar y) + \theta^-\psi_-(y,\bar y) + \theta^+\theta^-F(y,\bar y).
\tag{2.15}
\end{equation}
% [原文疑误: 原文 "The chirality constraint $\bar D_\pm\Phi=0$ implies that the superfield depends on the fermionic coordinates in via $y$, $\bar y$ and $\theta^\pm$" 一句中 "in via" 语意不通，疑为 "only via"（只经由）之误，且原句只提"费米坐标"却同时列出 $y,\bar y$（玻色坐标的组合）。译文按上下文取"只经由 $y$、$\bar y$ 与 $\theta^\pm$ 依赖各坐标"，未擅自补充物理内容。]

我们还要指出，扭手征超场在移位坐标 $\tilde y=z-i\theta^-\bar\theta^-$、$\bar y=\bar z-i\theta^+\bar\theta^+$ 下也有类似的展开：

\begin{equation}
\tilde\Phi(\tilde y,\bar y;\theta^\pm,\bar\theta^\pm) = \nu(\tilde y,\bar y) + \theta^+\bar\chi_+(\tilde y,\bar y) + \bar\theta^-\chi_-(\tilde y,\bar y) + \theta^+\bar\theta^-E(\tilde y,\bar y).
\tag{2.16}
\end{equation}
% [Tutor Note: 与 (2.15) 对照可见扭手征展开用的是 $\theta^+$ 与 $\bar\theta^-$ 这两个"混合"方向（而手征展开用 $\theta^+$ 与 $\theta^-$）。原因在于手征约束是 $\bar D_\pm\Phi=0$，扭手征约束是 $\bar D_+\tilde\Phi=D_-\tilde\Phi=0$；后者的可解方向正是 $\theta^+$、$\bar\theta^-$，故展开只含这两个 Grassmann 方向的单项式，$F$ 项也相应换成 $E$ 项。注意原文本式第二个玻色宗量写作 $\bar y$（无 tilde），与原书此处记号一致，照抄未改。]

% ==== 原文 PDF 第 14 页 ====

共轭超场有类似的分量展开。我们把手征超场 $\Phi^n$ 及其共轭 $\bar\Phi^{\bar n}$ 看作从超空间到某个复靶空间 $M$ 的映射（并复合以局部复坐标的选取）。扭变版本有类似的解释。

利用上述手征超场的分量表达式，可以直接计算超荷 $Q_\pm$、$\bar Q_\pm$ 的作用；扭手征超场及其共轭场可类似处理。首先，作用在手征超场上时，即在坐标 $(y,\bar y;\theta^\pm,\bar\theta^\pm)$ 中，矢量场 $Q_\pm$、$\bar Q_\pm$ 如下给出：

\begin{equation}
\begin{aligned}
Q_+ &= \partial_{\theta^+} & \bar Q_+ &= -\partial_{\bar\theta^+} - 2i\theta^+\partial_{\bar z}\\
Q_- &= \partial_{\theta^-} & \bar Q_- &= -\partial_{\bar\theta^-} + 2i\theta^-\partial_z
\end{aligned}\tag{2.17}
\end{equation}

由于手征超场只通过 $y$、$\bar y$ 依赖于 $\bar\theta^\pm$，我们可以放心地忽略 $\bar Q_\pm$ 的第一项。然后通过下式定义 $Q_\pm$ 对分量 $\phi,\psi_\pm,F$ 的作用：

\begin{equation}
\begin{aligned}
Q_\pm\Phi &= \psi_\pm \pm \theta^\mp F\\
&:= (Q_\pm\phi) - \theta^+(Q_\pm\psi_+) - \theta^-(Q_\pm\psi_-) + \theta^+\theta^-(Q_\pm F)
\end{aligned}\tag{2.18}
\end{equation}

由此可得 $Q_\pm$ 对分量场的作用如下：

\begin{equation}
\begin{aligned}
Q_\pm\phi &= \psi_\pm\\
Q_\pm\psi_\pm &= 0 & Q_\pm\psi_\mp &= \mp F\\
Q_\pm F &= 0
\end{aligned}\tag{2.19}
\end{equation}
% [Tutor Note: (2.18) 第二行的负号来自超荷的 Grassmann 奇性：$Q_\pm$ 作为奇算子从左作用、越过 $\theta^+$、$\theta^-$ 时各贡献一个 Koszul 符号，故展开中出现两个负号。用 (2.19) 代回可核对：对 $Q_+$ 有 $-\theta^-(Q_+\psi_-)=-\theta^-(-F)=+\theta^-F$，其余项为零，于是 $Q_+\Phi=\psi_++\theta^-F$；对 $Q_-$ 同理得 $\psi_--\theta^+F$，与第一行的 $\psi_\pm\pm\theta^\mp F$ 一致。]

类似地，$\bar Q_\pm$ 对分量场的作用为：

\begin{equation}
\begin{aligned}
\bar Q_+\phi &= 0 & \bar Q_-\phi &= 0\\
\bar Q_+\psi_+ &= 2i\partial_{\bar z}\phi & \bar Q_-\psi_+ &= 0\\
\bar Q_+\psi_- &= 0 & \bar Q_-\psi_- &= -2i\partial_z\phi\\
\bar Q_+F &= -2i\partial_{\bar z}\psi_- & \bar Q_-F &= -2i\partial_z\psi_+
\end{aligned}\tag{2.20}
\end{equation}
% [Tutor Note: (2.19) 与 (2.20) 的导数项（含系数 2）来自链式法则与显式项两部分的相加：$\bar Q_\pm$ 中的 $\partial_{\bar\theta^\pm}$ 通过 $\bar y$ 的 $\theta$ 依赖转为 $\partial_{\bar y}$，再与 (2.17) 中显式的 $\pm 2i\theta^\pm\partial_z$（或 $\partial_{\bar z}$）项合并，系数恰为 2。核对时须与 $y=z+i\theta^-\bar\theta^-$、$\bar y=\bar z-i\theta^+\bar\theta^+$ 的约定以及 $P_z=-i\partial_z$、$P_{\bar z}=-i\partial_{\bar z}$ 一并验算，否则 $\partial_z$ 与 $\partial_{\bar z}$ 的方向极易互换。]

超荷 $Q_\pm$、$\bar Q_\pm$ 对扭手征超场 $\tilde\Phi$ 的作用可以用类似方式得到。一个巧妙的确定办法是利用上一节描述的镜像自同构：镜像自同构自然地通过交换 $\theta^-\leftrightarrow\bar\theta^-$ 作用在超空间的奇坐标上。特别地，我们看到手征多重态被变换为扭手征多重态，其对应关系为 $\nu=\phi$、$\bar\chi_-=\psi_+$、$\chi_-=\psi_-$、$E=F$。于是，$Q_+$ 与 $\bar Q_-$（相应地 $\bar Q_+$ 与 $Q_-$）对扭手征多重态各分量的作用，由式 (2.19)（相应地式 (2.20)）在此对应关系下给出。
% [原文疑误: 对照 (2.16)，扭手征超场展开中的 $\theta^+$ 伙伴是 $\bar\chi_+$（即 (2.16) 写作 $\theta^+\bar\chi_+$），而此处镜像对应关系却写作 $\bar\chi_-=\psi_+$，下标与 (2.16) 不一致。按镜像交换 $\theta^-\leftrightarrow\bar\theta^-$：手征的 $\theta^+\psi_+$ 对应扭手征的 $\theta^+\bar\chi_+$，故应为 $\bar\chi_+=\psi_+$；原文的 $-$ 疑为下标笔误。译文照抄原文 $\bar\chi_-=\psi_+$，请读者以 (2.16) 为准。]

从这一初步分析中，我们在超多重态层面找到了镜像对称的一个实例：手征多重态镜像于扭手征多重态。诚然，这一镜像对称并不比上文第 2.2.1 节所描述的镜像对称深刻多少——扭手征多重态本质上就是被定义为手征多重态的镜像。更一般地，给定由手征多重态与矢量多重态构成的理论，就有一个由相同数据定义的、平凡地互为镜像的扭手征多重态与扭矢量多重态的理论。然而，镜像对称的大多数陈述要有趣得多：它们例如把手征多重态（及矢量多重态）的理论与另一个手征多重态（可能不含矢量多重态）的理论相交换！
